\section{Casualit\`a scalare e sottospazi invarianti del propagatore non vincolato}
\label{app:scalar-randomness}

Una scelta implementativa rilevante \`e l'uso di due numeri casuali scalari per particella e aggiornamento, riutilizzati in tutte le coordinate. Questa convenzione coincide con il codice in \texttt{optimizer\_internal.hpp} e nel backend asincrono. La sua conseguenza geometrica riguarda il tratto di dinamica privo di correzioni al dominio e va distinta da un'asserzione di convergenza.

\begin{proposition}[Invarianza dello span per la dinamica non corretta]
Siano $\bm x_i^0,\bm v_i^0\in\R^D$, $i=1,\ldots,N$, e
\begin{equation}
\mathcal S_0=\operatorname{span}\{\bm x_1^0,\ldots,\bm x_N^0,
\bm v_1^0,\ldots,\bm v_N^0\}.
\end{equation}
Supponiamo che i coefficienti del propagatore locale siano scalari comuni alle coordinate di una particella; che i leader sociali e i migliori personali siano posizioni gi\`a generate; e che nessuna correzione al dominio o operazione esterna introduca nuovi vettori. Allora, per ogni particella e ogni aggiornamento, posizioni, velocit\`a, migliori personali e leader appartengono a $\mathcal S_0$, con
$\dim\mathcal S_0\le\min\{D,2N\}$.
\end{proposition}

\begin{proof}
La tesi vale all'inizializzazione. Supponiamo valga al passo $k$. Il migliore personale e il leader sociale sono elementi di $\mathcal S_0$ per ipotesi induttiva. Se $q=a+b>0$, l'equilibrio $\bm x_{\rm eq}=(a\bm p+b\bm s)/q$ appartiene a $\mathcal S_0$. La soluzione di un sistema lineare scalare identico lungo ciascuna coordinata \`e una combinazione lineare di $\bm x^k$, $\bm v^k$ e $\bm x_{\rm eq}$, con gli stessi coefficienti scalari in tutte le coordinate; quindi $\bm x^{k+1},\bm v^{k+1}\in\mathcal S_0$. Per $q=0$ la soluzione \`e combinazione dei soli $\bm x^k$ e $\bm v^k$. La selezione di un nuovo best da posizioni generate preserva lo span. L'induzione conclude la prova e la stima dimensionale segue dal numero dei generatori.
\end{proof}

\paragraph{Limiti della proposizione.}
Il clipping componente per componente non \`e in generale una trasformazione lineare comune a tutte le coordinate, e pu\`o rompere l'invarianza. Anche l'introduzione di un immigrato proveniente da un'altra isola pu\`o ampliare lo span locale. Pertanto non \`e corretto interpretare $\dim\mathcal S_0\le2N$ come una restrizione assoluta dell'implementazione con vincoli attivi.

\paragraph{Interpretazione rispetto ai problemi studiati.}
Per $N=64$ il limite superiore $2N=128$ non impone uno span proprio nelle dimensioni interne $D=30,100$. Per il confronto esterno con $N=40$, il valore $2N=80$ \`e invece inferiore alle dimensioni di CEC 2013 LSGO e ad alcune dimensioni del protocollo esteso BBOB-LargeScale. Questa osservazione segnala un possibile meccanismo geometrico da misurare, non una prova di scarsa prestazione: il rango effettivo dipende anche da inizializzazione, box, traiettorie e correzioni. Nelle isole omogenee di dimensione locale $n$, il corrispondente limite \`e $2n$ prima di migrazione e correzioni. Un'eventuale versione con coefficienti casuali indipendenti per coordinata rappresenterebbe un algoritmo distinto e non viene confusa con la CPSO attualmente confrontata.
